% ===
\section{Problem Context}
% ===

% ------------------
%%% DENSIDADE INSTITUCIONAL E PERSISTÊNCIA DA DEGRADAÇÃO
% ------------------
%% Arquitetura de governança existente
The production-preservation problem unfolds in a context characterized not by institutional absence, but by an increasingly dense architecture of environmental governance. Over the past three decades, international agreements, national regulations, carbon-pricing systems, payments for environmental services, financial instruments, certification schemes, and supply-chain commitments have sought to constrain degradation or redirect economic activity toward less environmentally damaging outcomes. These mechanisms differ in legal status, object, territorial reach, and theory of change. Some establish aggregate objectives; others prohibit specific conduct, price measurable externalities, compensate observable conservation, condition access to capital, or disclose environmental attributes to consumers and investors. Their coexistence therefore does not constitute a single governance system. It constitutes a policy mix whose components intervene at different margins of the production-preservation problem.

%% Persistência e heterogeneidade da degradação
This institutional expansion has occurred alongside environmental trends that remain inconsistent with stated collective objectives. The synthesis of the Sixth Assessment Report documents widespread climate-related impacts and a narrowing set of feasible pathways for limiting additional warming \cite{ipcc2023synthesis}. The global biodiversity assessment, in turn, estimates that around one million species are threatened with extinction and that human activity has significantly altered approximately three quarters of the terrestrial environment \cite{ipbes2019global}. More recently, the \textit{Emissions Gap Report 2025} estimated that full implementation of the available nationally determined contributions would still imply warming of approximately $2.3$--$2.5^{\circ}\mathrm{C}$ over this century, while policies then in place implied approximately $2.8^{\circ}\mathrm{C}$ \cite{unep2025emissions}. In Brazil, official monitoring reported that approximately 729,000 km$^{2}$, or 17\% of the Amazon biome, had already been deforested by 2024 \cite{inpe2024monitoring}. These aggregate indicators coexist with periods of substantial improvement, including the sharp decline in Brazilian Amazon deforestation after 2004, which the literature associates with a combination of enforcement, protected areas, credit restrictions, and interventions in beef and soy supply chains \cite{nepstad2014slowing}. Environmental deterioration is therefore persistent and heterogeneous, but neither uniform nor institutionally invariant.

%% Interpretação das evidências
The coexistence of environmental degradation and environmental policy does not, by itself, demonstrate that existing instruments are ineffective. Such an inference would ignore the relevant counterfactual: degradation may have been greater in their absence, and individual instruments may produce substantial effects even when aggregate outcomes remain insufficient relative to collective goals. It is equally important not to collapse distinct evaluative concepts. Resource conversion may be privately rational when its expected private return exceeds the return to preservation; socially inefficient when relevant external costs remain unaccounted for; ecologically unsustainable when it erodes the functions required to maintain the system through time; or illegal when it violates an enforceable rule. These categories may overlap, but they are not equivalent. The problem examined here is consequently not reducible to misconduct, deficient ecological awareness, or an inherent moral opposition between private agents and collective welfare. It concerns the conditions under which individually rational decisions, interacting through markets and institutions, can generate outcomes that are collectively costly and ecologically degrading.

% ------------------
%%% INSTRUMENTOS LEGAIS E INSTITUCIONAIS
% ------------------
%% Coordenação internacional
International environmental agreements address part of this problem by establishing common objectives, reporting procedures, and mechanisms for cooperation. The Kyoto Protocol imposed quantified mitigation commitments primarily on developed economies and introduced emissions trading, joint implementation, and the Clean Development Mechanism \cite{unfccc1997kyoto}. The Paris Agreement broadened participation through nationally determined contributions, successive cycles of ambition, transparency arrangements, and provisions for finance, technology, and capacity building \cite{unfccc2015paris}. The Convention on Biological Diversity established a parallel framework centered on conservation, the sustainable use of biodiversity, and the fair and equitable sharing of benefits derived from genetic resources \cite{cbd1992convention}. These agreements are institutionally significant because climate stability and biodiversity conservation cannot be supplied through isolated national action alone. At the same time, their effectiveness depends on the ambition of collective commitments and on their translation into domestic law, budgets, administrative capacity, and credible implementation. International coordination can define objectives and reduce informational and strategic barriers among states, but it does not directly determine the land-use or investment decision of each productive agent.

%% Legislação brasileira
At the domestic level, command-and-control regulation translates environmental objectives into legally defined duties, standards, authorizations, and sanctions. The Brazilian Forest Code illustrates that this architecture is more heterogeneous than a simple opposition between prohibition and production. Law No. 12,651/2012 regulates Permanent Preservation Areas and Legal Reserves, establishes the Rural Environmental Registry and environmental regularization procedures, and provides for economic and financial mechanisms intended to support conservation and restoration \cite{brasil2012florestal}. The law therefore combines property-specific restrictions, information and monitoring infrastructure, regularization mechanisms, and potential incentives. Scientific assessments of the reform similarly identify both reductions in restoration obligations and new opportunities for monitoring, compensation, and conservation \cite{soaresfilho2014cracking,azevedo2017limits}. Characterizing the Forest Code exclusively as a rigid constraint would consequently overlook its hybrid design, just as characterizing it as an integrated solution would overlook the institutional conditions required to make its components operational.

%% Enforcement, compliance e custos privados
Regulation can alter expected returns by reducing the legal set of available actions and by attaching an expected sanction to non-compliance. Its practical effect, however, depends on the probability of detection, the speed and credibility of enforcement, the clarity of obligations, and the capacity of regulated agents to comply. These conditions also shape distributional incidence. A conservation rule may be socially justified while imposing concentrated opportunity and implementation costs on landholders whose private returns from land conversion are more immediate than the benefits they can appropriate from maintained ecological functionality. This concentration of costs does not by itself establish injustice or inefficiency, but it can increase resistance and reduce compliance when obligations are uncertain, monitoring is inconsistent, or compensatory mechanisms are unavailable. In the Brazilian Amazon, simulations of combined enforcement and conservation payments reveal precisely this trade-off: adding payments may reduce landholder income losses, but it raises implementation costs and its cost-effectiveness varies spatially with deforestation risk, opportunity costs, and enforcement capacity \cite{borner2015mixing}.

The distinction between formal rules and implemented rules is therefore essential. More than a decade after the 2012 reform, implementation of the Forest Code remains heterogeneous across Brazilian states because general legal provisions must be translated, monitored, and enforced through multiple institutional layers. Recent institutional analysis identifies incomplete validation of self-declared information in the Rural Environmental Registry and uneven specification of state-level regularization procedures as central operational bottlenecks \cite{oliveira2025implementation}. Earlier evidence likewise showed that a land registry could lower monitoring costs without necessarily eliminating illegal deforestation when validation, enforcement, and compliance incentives remained incomplete \cite{azevedo2017limits}. These findings do not imply that law is dispensable. They show that legal validity and practical effectiveness are different properties: the former defines obligations, while the latter depends on the institutions that connect those obligations to observable behavior and expected private payoffs.

% ------------------
%%% INTRUMENTOS ECONÔMICOS E FINANCEIROS
% ------------------
%% Precificação e licenças negociáveis
Economic instruments intervene more directly in relative prices. Environmental taxes increase the private cost of polluting activities, while cap-and-trade systems constrain aggregate emissions and allow regulated agents to exchange allowances. In principle, both can induce abatement where its marginal private cost is lowest. In practice, instrument choice also involves distributional effects, uncertainty, administrative requirements, innovation incentives, and political feasibility; no instrument is unconditionally superior across all relevant criteria \cite{goulder2008instrument}. Empirical evidence should therefore be interpreted with similar care. The European Union Emissions Trading System reduced emissions relative to its estimated counterfactual despite initially low allowance prices \cite{bayer2020euets}, and a systematic review of 80 ex-post evaluations found statistically significant reductions for at least 17 of the 21 carbon-pricing policies examined \cite{dobbeling2024carbon}. Carbon pricing is thus neither generally ineffective nor sufficient by definition.

Its limitation for the present problem is primarily one of scope. A carbon price assigns a monetary consequence to an additional unit of measured greenhouse-gas emissions or to the use of a regulated allowance. It does not automatically represent biodiversity, hydrological regulation, soil integrity, ecological thresholds, or the productive feedbacks generated by their joint condition. Even within climate policy, a price may be too narrow to induce infrastructure, innovation, and institutional changes whose benefits are uncertain, complementary, or realized only over long horizons \cite{rosenbloom2020pricing}. The relevant conclusion is not that pricing should be abandoned, but that an instrument designed for a measurable externality cannot be presumed to govern the entire ecological state on which production and welfare depend.

%% Pagamentos por serviços ambientais (PSA)
Payments for environmental services approach the incentive problem from the opposite direction by remunerating an agreed conservation action or environmental outcome. They can make preservation more attractive when payments are conditional, credible, and at least partly compensate opportunity and implementation costs. Reviews of the empirical literature find that such programs can generate additional conservation, while also emphasizing substantial variation in targeting, participation, monitoring, and estimated effects \cite{borner2017effectiveness,wunder2020payments}. Their performance depends on whether the contracted outcome is additional to what would have occurred without payment, whether degradation is displaced elsewhere, whether environmental gains persist after the contract, and whether the relevant service can be measured at reasonable cost. Moreover, compensation for a specified service does not necessarily incorporate the broader contribution of ecological functionality to production. A payment for avoided carbon emissions, for example, may preserve carbon while remaining spatially indifferent to hydrological connectivity or biodiversity value unless these dimensions are explicitly included in program design.

%% Instrumentos financeiros
Financial instruments operate at a different stage of the decision process. Green bonds, sustainability-linked finance, preferential credit, guarantees, and climate-related banking policies can influence the cost and availability of capital for environmentally designated activities. This channel matters because a profitable low-carbon or conservation investment may still fail to occur when it faces high upfront costs, long payback periods, credit constraints, or risks that private lenders cannot readily evaluate. For this reason, the financing problem cannot always be resolved by carbon pricing alone \cite{campiglio2016banking}. Evidence on corporate green bonds indicates that labelled issuance can attract investor demand and is associated with subsequent environmental commitments and performance improvements, particularly when external certification strengthens credibility \cite{flammer2021green}. Yet a financing label identifies the use or conditions of capital rather than the ecological state of a territory. Its environmental integrity still depends on taxonomies, disclosure, additionality, verification, and the substantive quality of the financed activity.

%% Governança privada
Private and collaborative governance adds monitoring, information, reputation, and market-access mechanisms to this policy mix. Certification programs, zero-deforestation commitments, non-governmental monitoring, and supply-chain agreements can transmit environmental requirements beyond the direct reach of public enforcement. Evidence concerning their role in the decline in Amazon deforestation indicates that private coordination can complement public policy \cite{nepstad2014slowing}. Nevertheless, voluntary commitments tend to be selective in membership, commodity coverage, and territorial reach, while reputational sanctions depend on traceability and on the value that firms and consumers attach to compliance. Collaboration can improve information and collective action, but its effectiveness depends on the structure of relationships among actors and on whether those relationships match the ecological and institutional problem being addressed \cite{bodin2017collaborative}. NGOs and private standards are therefore neither substitutes for the state nor merely symbolic actors; they are complementary governance mechanisms whose authority and capacity remain issue-specific.

% ------------------
%%% LIMITAÇÕES DOS INSTRUMENTOS EXISTENTES
% ------------------
%% Interação entre políticas
Because these instruments address different margins, environmental governance is necessarily organized as a policy mix. Regulation can establish a binding floor, pricing can reveal heterogeneous abatement costs, payments can compensate conservation, finance can relax investment constraints, and private governance can improve traceability or market discipline. Their combination may be complementary, but it may also create overlap, contradiction, regulatory arbitrage, or unstable expectations. The effectiveness of a policy mix depends not only on the presence of individual instruments, but also on the consistency of its objectives, the coherence of its components, the credibility of implementation, and the sequencing of change \cite{rogge2016policy}. Environmental policy under complex and uncertain conditions therefore requires portfolios of instruments and institutions rather than reliance on a universally dominant mechanism \cite{sterner2019policy}. Institutional density should not be confused with institutional integration.

%% Descompasso de objeto, escala, tempo e implementação
Four recurrent mismatches remain particularly relevant to the analytical problem developed in this article. The first concerns the \textit{object} of governance. Existing instruments generally regulate an emission, a prohibited act, a contractual service, an eligible investment, or a disclosure obligation. Ecological functionality, however, is multidimensional and jointly affected by production, preservation, spatial interaction, and ecological feedback. Governing one observable component may improve that component without determining the condition of the wider ecological system.

The second mismatch is \textit{spatial}. Property boundaries, municipal jurisdictions, national commitments, and spatially fungible markets do not necessarily coincide with watersheds, habitats, ecological corridors, climatic feedbacks, or other functional ecological units. Scale mismatches arise when responsibility is allocated at a level that does not correspond to the environmental process being governed, producing information loss, spillovers, and incomplete incentives \cite{cumming2006scale}. The literature on institutional fit similarly shows that sustainability depends on the correspondence among ecological processes, social organization, and governance arrangements rather than on decentralization or centralization in the abstract \cite{epstein2015fit}. Local governance can improve contextual information and connect benefits to affected agents, but cross-boundary effects still require coordination at broader levels.

The third mismatch is \textit{temporal and distributive}. Returns from resource conversion are often immediate and privately appropriable, whereas benefits from ecological maintenance may be delayed, uncertain, dispersed among beneficiaries, or reflected in avoided future losses. Legal obligations, payments, and finance can change parts of this comparison, but they do not automatically make the productive contribution of ecological functionality enter the agent's expected payoff. The fourth mismatch is \textit{institutional}: formal objectives must pass through rules, budgets, information systems, administrative organizations, and enforcement practices before they affect behavior. Where translation, monitoring, or policy credibility is weak, the incentive perceived by the agent can differ substantially from the incentive intended by the rule.

Taken together, the evidence does not imply that existing environmental instruments are generally ineffective. Command-and-control regulation can deter illegal conversion; carbon pricing can reduce emissions; payments can alter land-use decisions; financial instruments can redirect capital; and collaborative governance can improve monitoring and coordination. Their performance, however, remains conditional on measurement, targeting, enforcement, institutional capacity, policy credibility, and fit with the relevant ecological process. More importantly for the present analysis, these mechanisms typically govern specific externalities, transactions, or legal obligations. Their coexistence does not by itself establish an integrated relationship in which ecological functionality jointly determines production, enters private payoffs, and is governed at the spatial scale of the ecological processes involved.

This residual disconnect defines the transition from context to theory. The analytical gap is not the absence of environmental concern from economic policy, nor does it justify rejecting regulation, markets, finance, or collective governance. It is the absence, for the purposes of this article, of a unified representation connecting the ecological state, productive outcomes, private incentives, and territorially aligned institutions. Section 3 therefore examines how the theories of externalities, property rights, common-pool resources, and natural capital illuminate separate components of this problem, before Section 4 reconstructs their connection through the proposed framework of economic commensalism.

% end %
